\begin{abstract}
We propose that the obstruction to holography in de~Sitter space is \emph{epistemological}---a matter of
representation---rather than \emph{ontological}---a matter of which vacuum is realised. We trace it to the
non-paradoxical self-reference of binary language, in the precise sense of the Lawvere fixed-point theorem and
the Yanofsky diagonal $g(t)=\alpha(f(t,t))$, compounded by the dependence of quantum mechanics, as originally
formulated, on Boolean logic. Resolving the obstruction at the intersection of set theory and intuitionism, we
build a geometric framework carried on a single torus generated by $\golden=2\cos(\pi/5)$, whose modulus
$|\Om|=\tfrac12$ we propose as the discrete microstate that low-dimensional holography could use. From this one
microstate---without any ensemble average---we recover the AdS/CFT correspondence and, through the golden tower
$z=\golden^{\sigma}$, an M-theory type duality web, level by level. This construction happens to be inspired by Yuri Manin's hypothesised use of the same noncommutative tori that M-theory uses, but to approach the $\F_1$ program for the Riemann hypothesis. The observer enters not as a selector of the
vacuum---the microstate is the unique self-consistent solution, so there is no landscape to choose from---but as
the projection $\Pi$ that participates by accumulating information: its quantum Fisher information, conjugate to
the tower entropy and bounded by the ultraviolet cuts, builds the bulk geometry and, through the Landauer cost
of each bit and the Jacobson--Verlinde passage from entropy to curvature, sources Einstein's equations. The
objectivity threshold $f_{\mathrm{crit}}=\tfrac12$ coincides with the modulus of the microstate, and
entanglement is geometry ($\mathrm{ER}=\mathrm{EPR}$-like) read from that modulus. The torus is proposed as the single microstate that low-dimensional holography and de~Sitter appear to lack: the unique self-consistent solution its independent origins converge on, into which the AdS/CFT and M-theory corners project. Three companion papers develop,
independently, the arithmetic, operator, and identity threads the construction invokes.
\end{abstract}

\noindent\textbf{Keywords:} Holography; de~Sitter; M-theory; string theory; field with one element ($\F_1$); $\lambda$-rings; golden ratio; moduli spaces; formal verification.

\tableofcontents

\section{Introduction}\label{sec:intro}

\subsection{The problem: the observer in de~Sitter}\label{sec:problem}

\subsubsection{The participatory observer}\label{sssec:participatory}
The problem of holography in de~Sitter space is, as Witten frames it in \emph{Quantum Gravity in De~Sitter Space}, that General Relativity ``cannot be quantized---unless it is embedded in a more complete theory''~\cite{Witten01dS}: a theory in which the observer is not introduced from outside but ``described by the theory''~\cite{CLPW22}.
In a cosmology with positive cosmological constant there is no asymptotic region to which an external observer may
retreat; the observer is part of the system it measures and must therefore model the system that contains it. This
is what makes the de~Sitter case structurally different from the anti-de~Sitter case, which has been the most
successful approach to finding a UV complete description of quantum gravity, yet.

\subsection{Proposed solutions and their open issues}\label{sec:proposals}
In this section we review some of the most important proposed solutions to the problem, together with the issues
that each one leaves unresolved.

\subsubsection{Strominger: dS/CFT and non-unitarity}
In 2001 Strominger proposed the dS/CFT correspondence~\cite{Strominger01}: quantum gravity in a $D$-dimensional
de~Sitter space is equivalent to a Euclidean conformal field theory living at future (or past) infinity,
conceptualised as a $(D{-}1)$-dimensional sphere. The difficulty is that the resulting theory is generically
non-unitary: the boundary conformal weights
\begin{equation}
h_{\pm}=\frac{D-1}{2}\pm\sqrt{\Bigl(\frac{D-1}{2}\Bigr)^{2}-m^{2}\ell^{2}}
\end{equation}
become complex once $m^{2}\ell^{2}>(\tfrac{D-1}{2})^{2}$, so probabilities do not sum to one and the dual theory is
mathematically elusive.

\subsubsection{Susskind: the static patch}
Susskind, one of the originators of the holographic principle~\cite{Susskind95}, addressed the problem from the
standpoint of a real observer: holography should be formulated inside the cosmological horizon of a static
observer, with all interior information encoded on that horizon at one bit per Planck area,
\begin{equation}
S_{\mathrm{hor}}=\frac{A_{\mathrm{hor}}}{4G_{N}} .
\end{equation}
In recent work this static patch has been connected to advanced quantum-mechanical models such as double-scaled
SYK; the holographic encoding itself, however, remains a postulate rather than a derived correspondence.

\subsubsection{Witten: finite dimensionality}
Also in 2001, Witten~\cite{Witten01dS} analysed the quantum properties of de~Sitter and argued that, because the
horizon has finite area, the Hilbert space of quantum gravity in de~Sitter must be strictly finite-dimensional,
\begin{equation}
\dim\mathcal{H}<\infty .
\end{equation}
This collides directly with traditional conformal field theories, which possess infinitely many states, and
imposes severe constraints on any holographic attempt; the observer, in particular, must be included in the
description rather than idealised away~\cite{CLPW22}.

\subsubsection{The ensemble and vacuum selection}\label{sssec:ensemble}
Despite these proposals the picture is, as Nico Cooper puts it, always the same~\cite{Cooper2026} (a gravitational theory at weak coupling in a
$d$-dimensional anti-de~Sitter interior is captured by a conformal theory at strong coupling on its
$(d{-}1)$-dimensional edge, the cleanest realisations being string backgrounds that flow to supergravity. In low
dimensions this breaks down: two-dimensional JT gravity has no single boundary dual, and the bulk is recovered
only as an average over an \emph{ensemble} of boundary theories~\cite{SSS19,Cooper2026}, even though quantum
gravity should carry its own discrete, individual microstates---much as an ensemble of Hamiltonians reproduces
Wigner's nuclear decay rates while any single nucleus is governed by one. This ensemble
problem connects directly to the observer problem in de~Sitter: the absence of a definite spatial boundary confines
the observer to an averaged statistical description within their cosmological horizon, breaking the notion of a
single pure quantum state and forcing an effective description that reconciles the finiteness of the horizon with
the theoretical existence of fundamental microstates). The most widely held response is to abandon uniqueness altogether: it merges the Everettian
many-worlds interpretation of quantum mechanics with anthropic vacuum selection over the string
landscape---an estimated $\sim\!10^{500}$ metastable de~Sitter
vacua~\cite{BoussoPolchinski00,Douglas03,Susskind03}. Our vacuum is then one branch among many,
populated by eternal inflation and fixed by anthropic selection rather than derived, the branches
of the wavefunction being identified with the landscape's vacua rather than with a single
microstate~\cite{BoussoSusskind11}.

% =====================================================================
\subsection{Hypothesis: self-reference as the common cause}\label{sec:hypothesis}

\subsubsection{A matter of mathematical representation}\label{sssec:mathrep}
We propose that the problem may be epistemological---a matter of mathematical representation---and not ontological---a matter of
vacuum selection in which we do not know which vacuum we inhabit. We propose that the problem is how to describe something without referring to that same something one is describing. Across many scientific and mathematical worlds we
encounter geometric structures that are part of a whole identical to themselves: fractals. They occur even in
literature, among the rhetorical figures---the \emph{mise en abyme}, for example. Self-reference in mathematics,
however, is a profoundly delicate matter, one that wreaks havoc in domains far beyond the de~Sitter universe we
address here. The way the whole relates to the part appears as a concrete obstruction across many domains: in set
theory, as the collection of all collections that do not contain themselves; in the foundations of arithmetic, as a
true proposition that asserts its own unprovability and so escapes every proof; in the theory of truth, as the
impossibility of defining truth within the language itself; in computation, as the impossibility of deciding
whether an arbitrary procedure---itself included---ever halts; and in engineering, as the feedback of a system that
must represent itself. We propose that the essence of the problem lies in how to describe the object without referring to the object, akin to the figure--ground problem in the graphic arts and painting. In essence, it is the Hilbert--P\'olya operator problem and the Weil Rosetta Stone, together with the development of the $\F_1$ program---a way of referring to something without referring to that something directly: in the case of $\F_1$, referring to the zeros of $\zeta$, or to the primes, without referring to the primes or to the zeros of $\zeta$.

\subsubsection{The Lawvere--Yanofsky obstruction}
Let us first examine how the problem of self-reference wreaks havoc across diverse domains. In de~Sitter holography the observer is forced to classify itself
by a two-valued distinction---inside or outside the horizon, bulk or boundary, yes or no---and it is this binary
self-reference, not self-reference in general, that becomes obstructed. The obstruction is made precise by the
Lawvere fixed-point theorem~\cite{lawvere69}: in a cartesian closed category a point-surjective map
$f\colon T\to Y^{T}$ forces every endomorphism of $Y$ to have a fixed point, so a single fixed-point-free
endomorphism forbids such an $f$. Yanofsky's diagonal~\cite{yanofsky03} packages the classical paradoxes through
\begin{equation}\label{eq:LY}
g(t)=\alpha\bigl(f(t,t)\bigr),\qquad \alpha(x)=1-x,
\end{equation}
where the fixed-point-free involution $\alpha$ on the binary register $\{0,1\}$ is precisely what makes the
\emph{binary} self-application $f(t,t)$ contradictory. Our hypothesis is that the de~Sitter obstruction is an
instance of this binary mechanism, compounded by the dependence of quantum mechanics---as it was originally
formulated---on Boolean logic.

\subsubsection{Stratification as prohibition}
Set theory under classical bivalent logic resolves self-reference by \emph{prohibiting} it through stratification.
The axiom of extensionality defines objects purely by their membership extension, and the axiom schema of
separation restricts comprehension so that no set is formed by a predicate quantifying over the totality of
sets---because the predicate $x\notin x$ applied to the class of all sets produces Russell's paradox. The
resolution is the cumulative hierarchy: every set has a rank $\alpha$, with
\begin{equation}
x\in u \;\Longrightarrow\; \mathrm{rank}(x)<\mathrm{rank}(u),
\end{equation}
and no set has a complement that is also a set, so the Boolean algebra of union and intersection is available
within each level while global complementation is missing. A morphism translating between two domains $A$ and $B$
would have to occupy some rank $\alpha$ while its arguments and values refer across levels---precisely the
cross-level self-application that stratification forbids. The barrier between universes is therefore not an
additional restriction imposed on top of the paradox-avoidance mechanism: it \emph{is} that mechanism.
Stratification solves self-reference by making cross-level morphisms illegal, and cross-level morphisms are exactly
what inter-domain translation requires.

\subsubsection{Boolean logic is not the only logic}
The preceding obstruction is specific to the classical, bivalent setting. Boolean logic, however, is not the only
logic that exists, and the resolution we shall pursue lies in replacing prohibition by mediation---a possibility
opened by the non-Boolean logics reviewed next.

% =====================================================================
\subsection{Mathematical antecedents: non-Boolean strategies for self-reference}\label{sec:math}

The classical, bivalent prohibition of \S\ref{sec:hypothesis} is not the only way to confront self-reference; the
literature documents several \emph{distinct} strategies, each with its own mechanism, and they should not be
conflated into a single route. Historically the constructivist and intuitionist tradition questioned classical
logic from outside: Kant's account of number as a construction in the pure intuition of time is a philosophical,
not a mathematical, antecedent, followed by Kronecker's finitism, Poincar\'e's pre-intuitionism, the constructivism
of Borel and Lebesgue, and Brouwer's rejection (1907) of the excluded middle for infinite domains. The mechanisms
that bear directly on the Lawvere--Yanofsky obstruction, each handling self-reference in a different way, are the
following.

\begin{enumerate}
\item \textbf{Type-theoretic stratification} (Russell~\cite{russell08}, Martin-L\"of~\cite{martinlof75}) proceeds by
\emph{prohibition by level}: forbidding $X\in X$ through a hierarchy of universes $U_0\subset U_1\subset U_2$ breaks
the self-referential cycle at the foundational level. It was the first tower strategy, but, as in
\S\ref{sec:hypothesis}, it does so by making cross-level morphisms illegal.
\item \textbf{Paraconsistent logic} (Priest~\cite{priest79}) proceeds by \emph{tolerance}: contradictions are
accepted as non-explosive, so a system may contain a contradiction without becoming trivial, and the diagonal need
not be excluded at all.
\item \textbf{Tarski's decidable geometry}~\cite{tarski52} proceeds by \emph{removing the capacity for
self-reference}: the first-order theory of real closed fields, which contains Euclidean geometry $\R^n$, is
complete and decidable precisely because it cannot internally define $\N$. Without G\"odel numbering no formula can
be assigned a code, and the Yanofsky diagonal $g(t)=\alpha(f(t,t))$ \emph{cannot form}.
\item \textbf{Intuitionism, Heyting algebras and the topos} (Brouwer, Heyting; Curry--Howard; Lawvere) proceed by
\emph{mediation}: the internal logic is intuitionistic, valued in a Heyting algebra, and the implication $A\to B$ is
realised as an internal object---the exponential $B^A$, governed by the currying adjunction
$\mathrm{Hom}(C\times A,\,B)\cong\mathrm{Hom}(C,\,B^A)$ that also underlies the Lawvere statement~\eqref{eq:LY}.
Self-reference is carried by the exponential rather than forbidden.
\end{enumerate}

What these strategies have in common is structural rather than logical: several are \emph{tower} constructions.
Russell's hierarchy of types $U_0\subset U_1\subset U_2$ and Martin-L\"of's hierarchy of universes both resolve
self-reference by ranking objects in an ascending tower of levels indexed by a single rule, and this stratified,
level-indexed form recurs throughout the constructivist tradition.

% =====================================================================
\subsection{Physical antecedents}\label{sec:phys}

\subsubsection{Throats and vacuum selection}
In holography there already exist constructions that resemble these strategies, but from what we like to call the
Tarskian or geometric perspective. In the near-horizon limit of a brane
stack the geometry develops an anti-de~Sitter \emph{throat}~\cite{Maldacena98} whose radial coordinate plays the
role of a renormalisation-group scale: in Poincar\'e coordinates
\begin{equation}
ds^{2}=\frac{L^{2}}{z^{2}}\bigl(-dt^{2}+d\vec{x}^{2}+dz^{2}\bigr),
\end{equation}
the boundary $z\to 0$ is the ultraviolet of the dual theory, with $E_{\mathrm{SYM}}=(L/z)\,E_{\mathrm{proper}}$. The
bulk interior is encoded on the boundary through the holographic bound
\begin{equation}
S=\frac{A}{4G_N},
\end{equation}
one bit per Planck area, matching the $N^{2}$ degrees of freedom of the dual gauge theory~\cite{SusskindWitten98}. In
a \emph{warped} throat this count runs with the radial coordinate: in the Klebanov--Strassler/Tseytlin
cascade~\cite{KS00} the effective rank grows logarithmically,
\begin{equation}
N_{\mathrm{eff}}(r)\sim g_s M^{2}\,\log\!\frac{r}{r_s},
\end{equation}
so the throat \emph{gains degrees of freedom---``gains bits''---as it grows} toward the ultraviolet, terminating in
the deep infrared in a confining gauge group with a mass gap through a cascade of Seiberg dualities. This already resembles the Russellian type-theoretic hierarchy; what is missing
is a \emph{logical} selection criterion. In the orthodox formulation the rule was a strict duality, one bulk
universe to one boundary theory; but in the simplest controllable case, $2$d~JT gravity, the bulk cannot be matched
to a single discrete boundary theory, and the mathematics requires instead a statistical ensemble of random
matrices~\cite{SSS19}, in tension with the discrete microstates expected of quantum gravity. Beyond the anthropic
principle---selection by one's own existence---there is no rule isolating a single conformal field theory; this is
the vacuum-selection problem, magnified in de~Sitter by the landscape of $\sim 10^{500}$ vacua~\cite{Susskind03}.

\subsubsection{M-theory from the superpoint}\label{sssec:superpoint}
A second, independently developed string/M-theoretic construction selects its content logically rather than anthropically. Applying the Lawvere development of
Hegel---the \emph{unity of opposites} formalized as an adjunction~\cite{lawvere96unity}, carried
into the cohesive $\infty$-topos in which the superpoint lives~\cite{Schreiber13cohesive}---Huerta and
Schreiber~\cite{huerta2018m} prove that the super-Minkowski spacetimes of string/M-theory arise as iterated
\emph{maximal invariant central extensions} of the superpoint $\R^{0|1}$---the simplest super-Minkowski spacetime,
with neither Lorentz nor spinorial structure. Forming these extensions consecutively discovers the spacetimes that
carry superstrings and culminates in the ten- and eleven-dimensional cases,
\begin{equation}
\R^{0|1}\;\longrightarrow\;\cdots\;\longrightarrow\;\R^{9,1|16}\;\longrightarrow\;\R^{10,1|32},
\end{equation}
each step singled out by super-Lie-algebra cohomology rather than postulated; the full brane content---the
D-branes, the M5-brane and their dualities---then follows from the induced higher central extensions of the
\emph{brane bouquet}. The dimensional structure is in this sense algebraically necessary, with no anthropic choice.
Crucially, the construction is binary at its root. Supergeometry is $\Z_2$-graded: it rests on the cyclic group
$\Z_2=\{0,1\}$ with addition modulo two ($1+1=0$), the most elementary binary structure, which fixes the sign rule
of the superalgebra and hence the (anti)commutation of the supercharges that the invariant cocycles encode. This
$\Z_2$ is precisely the binary register on which the involution $\alpha(x)=1-x$ of \S\ref{sec:hypothesis} acts. The
sense in which this is binary is worth distinguishing: a classical bit is Boolean, taking the value $0$ or $1$; a
qubit is quantum, a superposition of the two; the superpoint's binary is geometric, its $\Z_2$-grading
distinguishing the even, bosonic directions of ordinary geometry (the ``$0$'') from the odd, fermionic direction
that carries supersymmetry (the ``$1$''), the superpoint $\R^{0|1}$ being the minimal seed consisting of that single
odd direction. M-theory begins here because it is the smallest possible seed---a single relation between two
distinct kinds of coordinate---from whose iterated extensions the eleven-dimensional fabric is generated. So,
although its selection is logical rather than anthropic, the construction remains within the binary register and
there is no observer to be found.

% =====================================================================
\subsection{Physico-mathematical antecedents: Manin's $\F_1$ program}\label{sec:manin}

\subsubsection{Quantum tori and numbers as functions}
Between string theory and pure mathematics there is an intermediate, experimental line of research that sought to
use the quantum tori of M-theory, or the rings of Borger, to approach the Riemann hypothesis: the $\F_1$-geometry
program associated with Yuri Manin. The $\F_1$ (``field with one element'') program is the one in which a number
of mathematicians---among them Manin, Deninger, Soul\'e, Kurokawa, Connes and Consani, and Borger---have sought to
prove the Riemann hypothesis by transporting Weil's proof of the Riemann hypothesis \emph{for curves over finite
fields} to the arithmetic of $\Z$: were $\operatorname{Spec}\Z$ realisable as a ``curve'' over an absolute base
$\F_1$, Weil's geometric argument might carry over.

Manin's quantum-tori strand~\cite{manin04} uses the noncommutative tori that Connes, Douglas and
Schwarz~\cite{connes98} showed arise as compactifications of M-theory, placing M-theory
and noncommutative/\allowbreak$\F_1$-geometry in contact. Such a torus is the noncommutative space
$T^2_\theta$ whose algebra
$A_\theta$ is generated by two unitaries $U,V$ obeying
\begin{equation}
VU=e^{2\pi i\theta}\,UV,\qquad \theta\in\R\setminus\Q,
\end{equation}
the irrational angle $\theta$ encoding the real multiplication that Manin attaches to real quadratic fields.

This is particularly significant for the Riemann hypothesis because, as Elvang, Herderschee and
Morales~\cite{Elvang2026} showed, the odd zeta values appear in the Veneziano string amplitude: the Wilson
coefficients $a_{2k-1,0}=\alpha^{\prime\,2k+1}\zeta(2k+1)$ left free by the effective field theory are fixed, in
the unique stringy completion, to the Veneziano values~\cite{Veneziano}. The odd zeta values thus tie the two
frameworks---arithmetic and string-theoretic---together beyond Manin's program itself.

\subsubsection{A geometric route to the Riemann hypothesis}
In his other program, \emph{Numbers as Functions}~\cite{manin13}, Manin develops the conception of numbers as
functions. Its arithmetic-differential strand replaces the derivative by the Fermat quotient,
\begin{equation}
\delta_p(a)=\frac{a-a^{p}}{p},\qquad \psi_p(a)\equiv a^{p}\!\!\pmod p,\qquad \psi_p\circ\psi_q=\psi_{pq},
\end{equation}
the lift $\psi_p$ being the Frobenius of a $\Lambda$-ring~\cite{buium05}. Here Manin approaches the Riemann
hypothesis from a geometric perspective intended to transcend the self-reference that proving it entails: one
cannot refer to the primes, nor to any zero, without importing circularity. The geometric route is meant to
sidestep this by working with the space of the problem rather than its arithmetic statement.

\subsubsection{Riemann's $\zeta$, holography and the primes: beyond Manin's programs}\label{sssec:zeta-primes}
Both of Manin's frameworks remained unfinished before his death in January 2023; both still lack a \emph{dynamical
completion}, and both programs remain open.

Beyond Manin's work, however, further relations have been established between string theory, the spectrum of the primes, the Riemann hypothesis and the problem of holography in de~Sitter. Two of the most significant and best known bear directly on the problem: the odd zeta values appear both
in the perturbative QED contributions to the electron's anomalous magnetic moment---where $\zeta(3)$ and $\zeta(5)$
enter the three-loop coefficient~\cite{LaportaRemiddi96}---and in the Veneziano amplitudes already mentioned. A
more recent and striking relation is the one between black-hole interiors and the primes obtained by Hartnoll, Yang
and De~Clerck~\cite{HartnollYang25,DeClerckHartnollYang25}: near a spacelike (BKL) singularity the Wheeler--DeWitt
wavefunctions become automorphic forms whose $L$-functions, expanded over their Euler product, take the form of the
partition function of a free gas of oscillators labelled by the primes---a ``primon gas''~\cite{Julia90}---the
five-dimensional case requiring the complex (Gaussian and Eisenstein) primes. Finally, and no less significant, is
the relation between the GUE (Gaussian Unitary Ensemble) of random matrices---whose statistics Montgomery found to
govern the pair correlation of the zeta zeros, as recognised by Dyson and verified numerically by
Odlyzko~\cite{Montgomery73,Odlyzko87}---and the elusiveness of the P\'olya operator: the Hilbert--P\'olya
conjecture that the zeros form the spectrum of some self-adjoint Hamiltonian.

Read in the light of holography, three points stand out: (i) the Veneziano amplitudes generate the odd zeta values; (ii) AdS/CFT has been most successful precisely for black holes, which have been in turn related to the primes and GUE random matrices; and (iii) the Hermitian operator of the zeros, like the discrete microstate of low-dimensional AdS/CFT, is GUE---itself a statistical description used specifically when no Hamiltonian is known. Taken together,
these coincidences strongly suggest that AdS/CFT bears a closer relation to the Riemann hypothesis, at the
mathematical level, than it might at first appear; and that, given the relations above, the absence of a single
concrete microstate in both cases may share the same mathematical---rather than physical---cause. It is on this
that we base our contention that the conjecture of selection among the $\sim 10^{500}$ vacua is mistaken, and that
the conflict may be principally epistemological rather than physical.

% =====================================================================
\subsection{The framework}\label{sec:framework}

\subsubsection{A geometric construction}
Taking as precedent the frameworks of string theory, holography and AdS/CFT, together with Manin's framework and
his reading of Borger's, and proceeding from a mainly geometric perspective, we develop a framework that sits
halfway between set theory and intuitionism. It is intuitionist in the classical sense---in that it takes the
observation of nature as its starting point rather than purely abstract mathematical axioms---yet it does not
conflict with classification under set-theoretic equivalence; self-reference is thereby mediated at the
intersection of the two systems, rather than prohibited. Conceptually it is a self-similar, interdimensional fractal that contains, is part
of, and unifies---within a single lattice---the symmetries of $e^{i\pi}$, $\N$, $\E^3$, the primes, the
half-integers and $\C$, by means of an algebraic ring, working as a contemporary Vitruvian module that uses $\golden$, trigonometry and cyclotomic pentagons,
mounted on a torus.

\subsubsection{The ring on the torus}
We combine Euler's identity, the modulus of $\C$ and the Vitruvian modulus into a single ring (the unifying isomorphism is established independently in~\cite{Corr})---a
construction in which the four structural components (origin, magnitude as ratio to the module, angular direction,
and equivalence under magnitudes) coincide with the data of $e^{i\pi}+1=0$. These are carried by the ring
\begin{equation}
\Rpcf=\Z[\golden,\golden^{-1},\tfrac12],
\end{equation}
whose self-similarity is seeded by the fixed-point equation
\begin{equation}\label{eq:phisq}
\golden^{2}=\golden+1,\qquad \golden=\tfrac{1+\sqrt5}{2}=2\cos\tfrac{\pi}{5}=\zeta_{10}+\zeta_{10}^{-1},
\end{equation}
which extends the role of $i^{2}=-1$ from $\R\to\C$ to the further embedding $\C\to\E^{3}$ via $z=\golden y$. The
ring lives on a torus $T^2_{\PCF}$ that maintains a relation between $2$D and $3$D symmetry without selecting bulk
or boundary; in this reading the closed string of string theory is reinterpreted as the ring $\Rpcf$ living on the
torus---the extended identity. The same $i$ that seeds $\C$ also fixes the modulus from the $\pi$-side: through
$e^{i\pi}=-1$ and the Gaussian integral one has $\Gamma(\tfrac12)=\sqrt\pi$, hence the half-factorial
$(\tfrac12)!=\tfrac12\sqrt\pi$, so the common modulus $\mu=\tfrac12$ arises on the $i$/$\pi$ side exactly as on the
$\golden$ side (proved in \S\ref{ssec:origins}, Lemma~\ref{lem:gamma-half} and Theorem~\ref{thm:half-factorial};
the $\pi\leftrightarrow\golden$ bridge is Proposition~\ref{prop:pi-bridge}).

\subsubsection{Tarskian descent}
The mechanism allows a descent from arithmetic to pure geometry that we call Tarskian, after the decidability of
the first-order theory of real-closed fields: arithmetic statements are recast as geometric ones, and numbers are
treated as ratios---in particular as trigonometric ratios through $\golden=2\cos(\pi/5)$. The descent is what lets
the construction sidestep the direct self-reference that the arithmetic statement of the problem would require.
The descent yields a single geometric--arithmetic invariant---the unit of scale, one bit---that integrates the two
origins of the modulus into one state (made precise in \S\ref{ssec:microstate}, where the commutative diagram of
Theorem~\ref{thm:mu-diagram} shows the five routes to $\mu=\tfrac12$ coincide).

\subsubsection{The golden tower and the multiplicative spectrum}
On this geometric substrate the module is expanded as a tower of powers that are Russellian in form but geometric
rather than logical. The tower is anchored by the algebraic bridge
\begin{equation}\label{eq:bridge}
\golden^{\lambda_{\log}}=2,\qquad \lambda_{\log}=\frac{\log 2}{\log\golden}=\log_{\golden}2,
\end{equation}
so that $\golden$ generates the multiplicative spectrum that carries the binary base $2^{p}$. The same tower
carries the ternary base, $\golden^{\log_\golden 3}=3$, so its growth mediates between the binary register ($2$)
and the ternary register ($3$), with $\mu=\tfrac12$ as the geometric mean: $\golden^{\mu\lambda_{\log}}=\sqrt2$,
$\golden^{\mu\log_\golden 3}=\sqrt3$ (proved in \S\ref{ssec:origins}, Proposition~\ref{prop:mediation}; the
mediated Mersenne form is Proposition~\ref{prop:mersenne-med} in \S\ref{ssec:spectrum}).

\subsubsection{The selection of $\golden$: non-paradoxical self-reference}\label{sssec:selectphi}
The first reason for the choice of $\golden$ is a non-paradoxical self-referent relation: \eqref{eq:phisq} is the
unique positive solution of $x=1+1/x$, mapping onto itself without contradiction and maintaining consistency across
$1$D, $2$D and $3$D. This is the same relation that, as Pacioli's \emph{De~Divina Proportione}~\cite{pacioli1509}
records, made $\golden$ central in art, alongside other ratios.

\subsubsection{The selection of $\golden$: the empirical trace}
The second reason is empirical: $\golden$ recurs in wave functions of light, electricity and magnetism, and in many
lattices. Fibonacci numbers and the golden ratio appear across nearly all domains of science wherever
self-organisation or minimum-energy configurations are at play~\cite{Pletser2018}, including the inflation factor
and the multifractal tight-binding spectrum of the Fibonacci quasicrystal~\cite{Jagannathan2021}, and the gain
eigenvalue of four-mode parametric coupling in nonlinear photonic crystals~\cite{Jedrkiewicz2018}.

A more recent and striking appearance is gravitational. Among the extremal Kerr--Newman black holes whose mass,
charge and angular momentum are fixed multiples of the irreducible mass, Sonnino \emph{et al.} identify a
distinguished family whose coefficients depend solely on $\golden$~\cite{Sonnino2024}. Within it lies a unique
black hole whose physical parameters converge on $\golden$: analysed by differential geometry in $\E^3$, both the
scale parameter (twice the irreducible mass) and the Gaussian curvature of the horizon surface at the poles equal
$\golden$---the configuration of highest symmetry admissible in the Kerr--Newman metric, with a proposed bearing
on the dark-energy density parameter $\Omega_\Lambda$ of the Kerr--Newman--de~Sitter family. Read against the
black-hole/prime relation in \S\ref{sssec:zeta-primes}, where the Wheeler--DeWitt wavefunctions near a
five-dimensional spacelike singularity become automorphic $L$-functions whose Euler product is a primon gas, the
two results place black holes, $\golden$ and the primes on the same object from complementary sides: the horizon
geometry selects $\golden$, and the interior selects the primes.

\subsubsection{The selection of $\golden$: the classification of primes}
The third reason is arithmetic. The golden ratio is the fundamental unit of $\Q(\sqrt5)$, and the quadratic
character $\chi_5$ classifies every rational prime; in the cyclotomic-pentagonal arithmetic of Grisales Herrera
this yields a classification of the primes by $\golden$, realised below through the Dedekind zeta of $\Q(\sqrt5)$.

\subsubsection{The selection of $\golden$: extremal irrationality}
The fourth reason is that $\golden$ is the number least well approximated by rationals.
The same fixed-point relation $\golden=1+1/\golden$ is the continued fraction
$[1;1,1,\dots]$, all of whose partial quotients are the minimum $1$, so by Hurwitz's
theorem~\cite{Hurwitz1891} its rational approximations converge more slowly than those of
any other irrational: $\golden$ is maximally obstructed \emph{as a fraction}. Yet the same
number is exact \emph{as a ratio}---$\golden=2\cos(\pi/5)=\zeta_{10}+\zeta_{10}^{-1}$
\eqref{eq:phisq}, the diagonal-to-side ratio of the pentagon. The object whose arithmetic
representation is worst is the one whose geometric representation is exact, so it is
precisely $\golden$ that lets the construction carry the modulus as a ratio while never
resolving into a rational sublattice.

\subsubsection{The selection of $\golden$: optimal packing and the information bound}
The fifth reason joins the arithmetic to the physical. Extremal irrationality is, read
geometrically, \emph{optimal packing}: because $\golden$ resonates with no rational
sublattice, a distribution stepped by $\golden$ fills space with maximal uniformity---the
principle behind phyllotaxis, where seeds set at the golden angle leave no gaps. Read
informationally, the same property is a \emph{bound}: the most uniform packing is the
densest record of distinguishable states, so $\golden$ is the ratio at which a region holds
the most information it can hold. This we propose as a golden-ratio reading of the
holographic principle---a limiting quantity of information per region---and it is what makes
$\golden$ the natural carrier of the modulus $\mu=\tfrac12$, whose binary entropy
$H(\tfrac12)=1$ is exactly one bit per fundamental cell. The reading is realised
quantitatively in \S\ref{ssec:object}, where the modulus saturates the Bekenstein
bound~\eqref{eq:bridge-BH}.

\subsubsection{A single microstate}
The combination of the foregoing produces the mathematical equivalent of a \emph{single microstate}: a structure
describable by a Hermitian operator and manifesting as one specific torus carrying a string. In this paper the microstate enters geometrically, through that torus and its
invariants---which is all the derivations below require; its realization as a Hermitian operator is built,
from the same invariants, in the companion Hilbert--P\'olya
paper~\cite{HP}. Here we establish the mathematical framework from string theory and then trace how it connects to
physics and to the observer.

\subsection{The physical program}\label{sec:physics}
Sections~\ref{sec:derivations} and~\ref{sec:implications} demonstrate that the mathematical
apparatus assembled in Section~\ref{sec:methods}---the single microstate $|\Om|=\tfrac12$---makes
contact with physics. The aim is not to propose new physics, but to establish that this apparatus
meets criteria already fixed in the literature: each physical result is obtained by satisfying a
requirement that an accepted framework had left open.

\subsubsection{Established frameworks and their open criteria}\label{sssec:frameworks}
The physical content starting point is Maldacena's conjecture that
``M/string theory on various Anti-deSitter spacetimes is dual to various conformal field
theories''~\cite{Maldacena98}. We demonstrate how the mathematical framework allows the
holographic and the M-theoretic descriptions to meet in one object, of which the microstate is a
projection. Section~\ref{sec:derivations} presents
the microstate's relation to the 't~Hooft--Susskind holographic bound~\cite{Susskind95}, to the AdS
boundary conditions~\cite{Maldacena98}, and to the $\mathrm{ER}=\mathrm{EPR}$
bridge~\cite{MaldacenaSusskind13}; Section~\ref{sec:implications} presents its relation to Witten's
unitarity and entropy conjectures for de~Sitter space~\cite{Witten01dS} and to the observer
requirement of Chandrasekaran, Longo, Penington and Witten~\cite{CLPW22}. The equations are given
in the sections indicated and collected in the demonstrative lines of Sections~\ref{ssec:spine}
and~\ref{ssec:spine4}.

\subsubsection{Plan of the paper and division of labour}
The remainder of this paper proceeds as follows. Section~\ref{sec:methods} establishes the mathematical framework
and constructs the single microstate from its string-theoretic ingredients: the two geometric origins of
$\mu=\tfrac12$, the uniqueness of the arity that blocks the diagonal, the base ring on the torus and its golden
tower, and the convergence on $\mu$, closing with the formal verification. Section~\ref{sec:derivations} derives the
physics, obtaining all of AdS/CFT and the M-theory type duality web from that single microstate at each level.
Section~\ref{sec:implications} turns to the observer and to gravity: the observer is the projection $\Pi$, it does
not select a vacuum but accumulates information, and that accumulation yields, through the
Landauer--Jacobson--Verlinde chain, Einstein's equations. Section~\ref{sec:conclusions} concludes.
\begin{table}[H]\centering\small
\caption{Notation and key definitions.}\label{tab:notation}
\begin{tabular}{@{}ll@{}}
\hline
Symbol / object & Meaning \\
\hline
$\golden=2\cos(\pi/5)=\tfrac{1+\sqrt5}{2}$ & the golden ratio; generator of the torus \\
$\lng=\log2/\log\golden$ & golden logarithmic scale \\
$\mu=\tfrac12$ & the modulus, $\mu=|\Om|$ \\
$\sigma$ & tower level; $\sigma=\tfrac32$ at the spectral point \\
$\Om$ & the microstate; $|\Om|=\tfrac12$ \\
$T^2_{\PCF}=\C/\Lambda,\ \tau=i$ & the torus (geometric home) \\
$\Rpcf=\Z[\golden,\golden^{-1},\tfrac12]$ & the base ring \\
$z=\golden^{\sigma}$ & the golden tower / throat coordinate \\
$\Nmodes=\lfloor\pi\golden^{\sigma}\rfloor$ & mode count at level $\sigma$ \\
$\epsz$ & ultraviolet cutoff (certainty cell) \\
$\Mpcf$ & PCF mass scale \\
$\GN=\tfrac12$ & Newton constant \\
$c=3$ & central charge $=$ colour arity $=\lfloor\pi\rfloor=\chi(\mathbb{CP}^2)$ \\
$\SBH=A/4\GN$ & horizon (Bekenstein--Hawking) entropy \\
$f_{\mathrm{crit}}=\tfrac12$ & objectivity threshold \\
$\Pi:\E^3\to\C$ & the observer projection \\
\hline
\multicolumn{2}{@{}l}{\emph{Key definitions:}}\\
microstate & the single object $\Om$, $|\Om|=\tfrac12$, with no ensemble average \\
shared signature & $|\Om|=\tfrac12,\ \ell=1,\ \GN=\tfrac12,\ c=3,\ \mathrm{GKP}=\tfrac34$ \\
golden tower & the levels $z=\golden^{\sigma}$ of the construction \\
\hline
\end{tabular}
\end{table}


% =====================================================================


% ===================================================================

% =====================================================================

